% Unit 061 terjemahan Bahasa Indonesia dari chapter4.tex baris 2116--2267.
% Sumber: Methods of Algebra, Volume 2, commit
% 9a5803ff77dd3257484cb177f851a73770a59dd3 (CC BY 4.0).
% stable-unit-id: o014.aljabr2.chapter4.unbounded-derived-functors
% Cakupan: Bagian 4.11 lengkap; berhenti sebelum Bagian 4.12 pada baris 2268.

% segment-id: o014.li2.chapter4.unbounded-derived-functors.h001
\section{Funktor Turunan Tak Terbatas}\label{sec:unbounded-derived}
\hypertarget{o014.aljabr2.chapter4.unbounded-derived-functors}{ }
% segment-id: o014.li2.chapter4.unbounded-derived-functors.x001
\index{funktor turunan!tak terbatas (unbounded)}

% segment-id: o014.li2.chapter4.unbounded-derived-functors.p001
Tinjau funktor aditif $F:\mathcal{A}\to\mathcal{A}'$ antarkategori abelian.
Bagian ini kembali ke Definisi~\ref{def:derived-functor} dan berfokus pada
versi funktor turunan untuk kategori turunan tak terbatas,
% segment-id: o014.li2.chapter4.unbounded-derived-functors.q001
\[ \mathrm{R}F, \; \mathrm{L}F: \cate{D}(\mathcal{A}) \to \cate{D}(\mathcal{A}'). \]

% segment-id: o014.li2.chapter4.unbounded-derived-functors.p002
Untuk menjamin keberadaan $\mathrm{R}F$ atau $\mathrm{L}F$ dan melakukan
perhitungan dengannya, diperlukan konsep kompleks K-injektif atau kompleks
K-projektif dari Definisi~\ref{def:K-resolutions}, beserta resolusi
yang bersesuaian.

% segment-id: o014.li2.chapter4.unbounded-derived-functors.le001
\begin{lema}\label{prop:K-injective-zero}
	Misalkan $\mathcal{A}$ kategori abelian. Jika
	$Z\in\Obj(\cate{K}(\mathcal{A}))$ sekaligus merupakan kompleks K-injektif
	(atau K-projektif) dan asiklik, maka $Z=0$.
\end{lema}
% segment-id: o014.li2.chapter4.unbounded-derived-functors.pf001
\begin{bukti}
	Lema~\ref{prop:K-injective-criterion} menyiratkan
	$\Hom_{\cate{K}(\mathcal{A})}(Z,Z)=0$.
\end{bukti}

% segment-id: o014.li2.chapter4.unbounded-derived-functors.p003
Mula-mula kita nyatakan suatu hasil umum. Tetapkan kategori-kategori
bertriangulasi $\mathcal{D}_1$, $\mathcal{D}_2$, $\mathcal{D}'$ beserta
subkategori-subkategori bertriangulasi jenuhnya
$\mathcal{N}_1$, $\mathcal{N}_2$, $\mathcal{N}'$.

% segment-id: o014.li2.chapter4.unbounded-derived-functors.le002
\begin{lema}\label{prop:K-cats}
	Misalkan diberikan bifunktor bertriangulasi
	$G:\mathcal{D}_1\times\mathcal{D}_2\to\mathcal{D}'$, definisikan
	subkategori penuh $\mathcal{K}_2^G$ dari $\mathcal{D}_2$ melalui
% segment-id: o014.li2.chapter4.unbounded-derived-functors.q002
	\[ \Obj(\mathcal{K}_2^G) = \left\{ X_2 \in \Obj(\mathcal{D}_2) : X_1 \in \Obj(\mathcal{N}_1) \implies G(X_1, X_2) \in \Obj(\mathcal{N}') \right\}. \]
	Maka $\mathcal{K}_2^G$ merupakan subkategori bertriangulasi jenuh. Setelah
	menukar posisi kedua variabel, pernyataan yang bersesuaian tetap berlaku bagi
	$\mathcal{K}_1^G\subset\mathcal{D}_1$.
\end{lema}
% segment-id: o014.li2.chapter4.unbounded-derived-functors.pf002
\begin{bukti}
	Isomorfisme $G(X_1,X_2[1])\simeq G(X_1,X_2)[1]$ menunjukkan bahwa
	$\mathcal{K}_2^G$ tertutup terhadap pergeseran. Demikian pula, misalkan
	$I\to J\to K\xrightarrow{+1}$ segitiga terbedakan dengan
	$I,K\in\Obj(\mathcal{K}_2^G)$. Dari segitiga terbedakan
% segment-id: o014.li2.chapter4.unbounded-derived-functors.q003
	\[ G(X_1, I) \to G(X_1, J) \to G(X_1, K) \xrightarrow{+1} \]
	segera terlihat bahwa
	$X_1\in\Obj(\mathcal{N}_1)\implies
	G(X_1,J)\in\Obj(\mathcal{N}')$. Sifat jenuhnya sudah jelas.
\end{bukti}

% segment-id: o014.li2.chapter4.unbounded-derived-functors.p004
Hasil ini dapat diterapkan pada
$\Hom:\mathcal{A}^{\opp}\times\mathcal{A}\to\cate{Ab}$ dan bifunktor
bertriangulasi yang bersesuaian, yaitu $\cate{K}_{\Pi}\Hom$.
Contoh~\ref{eg:Hom-bicplx} telah menunjukkan bahwa bifunktor yang terakhir ini
pada dasarnya merupakan kompleks $\Hom$:
% segment-id: o014.li2.chapter4.unbounded-derived-functors.e001
\begin{equation}\label{eqn:CPiHom}\begin{gathered}
	(\cate{C}_{\Pi} \Hom)(\sigma^{-1} X_1, X_2) \simeq \Hom^\bullet\left( X_1, X_2 \right), \\
	(\cate{K}_\Pi \Hom)(\sigma^{-1} X_1, X_2) \simeq \Hom^\bullet\left( X_1, X_2 \right)\; \text{sebagai citra dalam $\cate{K}(\cate{Ab})$},
\end{gathered}\end{equation}
% segment-id: o014.li2.chapter4.unbounded-derived-functors.p005
Di sini
$\sigma:\cate{C}(\mathcal{A}^{\opp})\rightiso
\cate{C}(\mathcal{A})^{\opp}$ dan
$\cate{K}(\mathcal{A}^{\opp})\rightiso\cate{K}(\mathcal{A})^{\opp}$
adalah isomorfisme dari Definisi--Proposisi~\ref{def:sigma} dan
Proposisi~\ref{prop:sigma-homotopy}. Versinya untuk kategori turunan akan
diperlukan kemudian.

% segment-id: o014.li2.chapter4.unbounded-derived-functors.pr001
\begin{proposisi}\label{prop:K-injective-F-injective}
	Semua kompleks K-injektif (atau K-projektif) membentuk subkategori
	bertriangulasi jenuh $\mathcal{I}$ (atau $\mathcal{P}$) dari
	$\cate{K}(\mathcal{A})$. Jika $\mathcal{A}$ mempunyai cukup banyak kompleks
	K-injektif (atau K-projektif), maka untuk setiap funktor aditif
	$F:\mathcal{A}\to\mathcal{A}'$, subkategori $\mathcal{I}$ (atau
	$\mathcal{P}$) bersifat $F$-injektif (atau $F$-projektif).
\end{proposisi}
% segment-id: o014.li2.chapter4.unbounded-derived-functors.pf003
\begin{bukti}
	Cukup dibahas kasus K-injektif. Bagian pertama diperoleh dengan menerapkan
	Lema~\ref{prop:K-cats} pada $G=\cate{K}_{\Pi}\Hom$, sambil menggunakan
	Definisi~\ref{def:K-resolutions} tentang kompleks K-injektif dan
	\eqref{eqn:CPiHom}: subkategori $\mathcal{K}_2^G$ di sana tepat merupakan
	$\mathcal{I}$ di sini.

	Misalkan $\mathcal{A}$ mempunyai cukup banyak kompleks K-injektif. Jika $I$
	merupakan kompleks K-injektif yang asiklik, maka
	Lema~\ref{prop:K-injective-zero} menyiratkan $I=0$, sehingga
	$(\cate{K}F)I$ tentu saja asiklik. Dengan demikian, seluruh syarat dalam
	Definisi~\ref{def:F-injective} bagi subkategori bertriangulasi $F$-injektif
	terpenuhi.
\end{bukti}

% segment-id: o014.li2.chapter4.unbounded-derived-functors.p006
Tuliskan $Q:\cate{K}(\mathcal{A})\to\cate{D}(\mathcal{A})$ dan
$Q':\cate{K}(\mathcal{A}')\to\cate{D}(\mathcal{A}')$ untuk funktor-funktor
lokalisasi.

% segment-id: o014.li2.chapter4.unbounded-derived-functors.co001
\begin{korolari}
	Jika $\mathcal{A}$ mempunyai cukup banyak kompleks K-injektif, maka $F$
	mempunyai funktor turunan kanan
	$\mathrm{R}F:\cate{D}(\mathcal{A})\to\cate{D}(\mathcal{A}')$, sedemikian
	sehingga $\mathrm{R}F(QX)\simeq Q'\cate{K}F(X)$ apabila $X$ kompleks
	K-injektif.

	Secara dual, jika $\mathcal{A}$ mempunyai cukup banyak kompleks K-projektif,
	maka $F$ mempunyai funktor turunan kiri
	$\mathrm{L}F:\cate{D}(\mathcal{A})\to\cate{D}(\mathcal{A}')$, sedemikian
	sehingga $\mathrm{L}F(QX)\simeq Q'\cate{K}F(X)$ apabila $X$ kompleks
	K-projektif.
\end{korolari}
% segment-id: o014.li2.chapter4.unbounded-derived-functors.pf004
\begin{bukti}
	Terapkan hasil Proposisi~\ref{prop:K-injective-F-injective} pada kerangka umum
	Proposisi~\ref{prop:localization-injective}.
\end{bukti}

% segment-id: o014.li2.chapter4.unbounded-derived-functors.p007
Sekarang tetapkan kategori-kategori abelian $\mathcal{A}_1$, $\mathcal{A}_2$,
dan $\mathcal{B}$. Untuk bifunktor aditif
$F:\mathcal{A}_1\times\mathcal{A}_2\to\mathcal{B}$, tinjau versi tak
terbatas dari bifunktor turunan $\mathrm{R}F$ (atau $\mathrm{L}F$) dalam
Definisi~\ref{def:derived-bifunctor}. Selanjutnya selalu diasumsikan bahwa
$\mathcal{B}$ mempunyai produk terhitung (atau koproduk terhitung).

% segment-id: o014.li2.chapter4.unbounded-derived-functors.p008
Secara konkret, kita hendak menentukan $\mathrm{R}F$ (atau $\mathrm{L}F$)
menggunakan subkategori $\mathcal{I}_1$, $\mathcal{I}_2$ (atau
$\mathcal{P}_1$, $\mathcal{P}_2$) di atas.

% segment-id: o014.li2.chapter4.unbounded-derived-functors.le003
\begin{lema}\label{prop:bifunctor-via-K-injective-prep}
	Untuk $i=1,2$, ambil subkategori bertriangulasi $\mathcal{I}_i$ (atau
	$\mathcal{P}_i$) dari $\cate{K}(\mathcal{A}_i)$ seperti dalam
	Proposisi~\ref{prop:K-injective-F-injective}.
% segment-id: o014.li2.chapter4.unbounded-derived-functors.l001
	\begin{compactitem}
% segment-id: o014.li2.chapter4.unbounded-derived-functors.i001
		\item Jika $\mathcal{A}_1$ dan $\mathcal{A}_2$ mempunyai cukup banyak
		kompleks K-injektif, maka $\mathcal{I}_1\times\mathcal{I}_2$ bersifat
		$F$-injektif.
% segment-id: o014.li2.chapter4.unbounded-derived-functors.i002
		\item Jika $\mathcal{A}_1$ dan $\mathcal{A}_2$ mempunyai cukup banyak
		kompleks K-projektif, maka $\mathcal{P}_1\times\mathcal{P}_2$ bersifat
		$F$-projektif.
	\end{compactitem}
\end{lema}
% segment-id: o014.li2.chapter4.unbounded-derived-functors.pf005
\begin{bukti}
	Cukup dibahas kasus K-injektif. Kuncinya ialah menunjukkan bahwa jika
	$X_1\in\Obj(\mathcal{I}_1)$, maka
	$(\cate{K}_{\Pi}F)(X_1,\cdot)$ memetakan setiap objek asiklik $Z$ dari
	$\mathcal{I}_2$ ke objek asiklik dalam $\cate{K}(\mathcal{B})$. Namun,
	Lema~\ref{prop:K-injective-zero} memberikan $Z=0$, sedangkan
	$(\cate{K}_{\Pi}F)(X_1,\cdot)$ merupakan funktor aditif, sehingga pernyataan
	ini langsung berlaku. Hasil yang bersesuaian diperoleh setelah menukar indeks
	$1$ dan $2$.
\end{bukti}

% segment-id: o014.li2.chapter4.unbounded-derived-functors.p009
Tuliskan $Q:\cate{K}(\mathcal{B})\to\cate{D}(\mathcal{B})$ dan
$Q_i:\cate{K}(\mathcal{A}_i)\to\cate{D}(\mathcal{A}_i)$ untuk
funktor-funktor lokalisasi ($i=1,2$).

% segment-id: o014.li2.chapter4.unbounded-derived-functors.pr002
\begin{proposisi}\label{prop:bifunctor-via-K-injective}
	Untuk $F:\mathcal{A}_1\times\mathcal{A}_2\to\mathcal{B}$, misalkan
	$\mathcal{A}_1$ dan $\mathcal{A}_2$ sama-sama mempunyai cukup banyak
	kompleks K-injektif (atau K-projektif). Maka bifunktor turunan
	$\mathrm{R}F$ (atau $\mathrm{L}F$) ada. Jika
	$X_i\in\Obj(\cate{K}(\mathcal{A}))$ merupakan kompleks K-injektif (atau
	K-projektif) untuk $i=1,2$, maka terdapat isomorfisme kanonik
% segment-id: o014.li2.chapter4.unbounded-derived-functors.q004
	\[ \mathrm{R}F(Q_1 X_1, Q_2 X_2) \simeq Q (\cate{K}_{\Pi} F)(X_1, X_2) \quad \text{(atau $\mathrm{L}F(Q_1 X_1, Q_2 X_2) \simeq Q (\cate{K}_{\oplus} F)(X_1, X_2)$) }. \]
\end{proposisi}
% segment-id: o014.li2.chapter4.unbounded-derived-functors.pf006
\begin{bukti}
	Terapkan langsung hasil Lema~\ref{prop:bifunctor-via-K-injective-prep} pada
	Proposisi~\ref{prop:localization-injective-bifunctor}.
\end{bukti}

% segment-id: o014.li2.chapter4.unbounded-derived-functors.p010
Untuk bifunktor tertentu $F$, sering kali dapat dipilih subkategori
$F$-injektif (atau $F$-projektif) yang lebih besar daripada subkategori
kompleks K-injektif (atau K-projektif). Funktor $\RHom$ yang akan dibahas
berikut merupakan salah satu contohnya.

% segment-id: o014.li2.chapter4.unbounded-derived-functors.th001
\begin{teorema}[$\RHom$ tak terbatas]\label{prop:unbounded-RHom}
% segment-id: o014.li2.chapter4.unbounded-derived-functors.x002
	\index{funktor $\RHom$}
% segment-id: o014.li2.chapter4.unbounded-derived-functors.x003
	\index[sym1]{RHom@$\RHom$}
	Misalkan kategori abelian $\mathcal{A}$ mempunyai cukup banyak kompleks
	K-injektif atau cukup banyak kompleks K-projektif. Untuk menyederhanakan
	notasi, tuliskan semua funktor lokalisasi sebagai $Q$.
% segment-id: o014.li2.chapter4.unbounded-derived-functors.l002
	\begin{enumerate}[(i)]
% segment-id: o014.li2.chapter4.unbounded-derived-functors.i003
		\item Bifunktor
		$\Hom=\Hom_{\mathcal{A}}:\mathcal{A}^{\opp}\times\mathcal{A}
		\to\cate{Ab}$ mempunyai bifunktor turunan kanan
		$\cate{D}(\mathcal{A}^{\opp})\times\cate{D}(\mathcal{A})
		\to\cate{D}(\cate{Ab})$. Melalui ekuivalensi $\sigma$ dari
		Proposisi~\ref{prop:derived-cat-op}, bifunktor itu juga dapat ditulis sebagai
% segment-id: o014.li2.chapter4.unbounded-derived-functors.q005
		\[ \RHom: \cate{D}(\mathcal{A})^{\opp} \times \cate{D}(\mathcal{A}) \to \cate{D}(\cate{Ab}). \]
		Bahkan, jika $\mathcal{A}$ mempunyai cukup banyak kompleks K-injektif
		(atau K-projektif), maka
		$\cate{K}(\mathcal{A})^{\opp}\times\mathcal{I}$ (atau
		$\mathcal{P}^{\opp}\times\cate{K}(\mathcal{A})$) bersifat
		$\Hom$-injektif.
% segment-id: o014.li2.chapter4.unbounded-derived-functors.i004
		\item Misalkan $\mathcal{A}$ mempunyai cukup banyak kompleks K-injektif
		dan $X_2\to I_2$ merupakan resolusi K-injektif. Resolusi ini menginduksi
% segment-id: o014.li2.chapter4.unbounded-derived-functors.q006
		\[ \RHom(Q X_1, Q X_2) \rightiso Q \Hom^\bullet(X_1, I_2). \]
		Sebagai kasus khusus,
		$\RHom(QX_1,\cdot)\simeq
		\mathrm{R}\left(\Hom(X_1,\cdot)\right)$.
% segment-id: o014.li2.chapter4.unbounded-derived-functors.i005
		\item Misalkan $\mathcal{A}$ mempunyai cukup banyak kompleks K-projektif
		dan $P_1\to X_1$ merupakan resolusi K-projektif. Resolusi ini menginduksi
% segment-id: o014.li2.chapter4.unbounded-derived-functors.q007
		\[ \RHom(Q X_1, Q X_2) \rightiso Q \Hom^\bullet(P_1, X_2). \]
		Sebagai kasus khusus,
		$\RHom(\cdot,QX_2)\simeq
		\mathrm{R}\left(\Hom(\cdot,X_2)\right)$.
% segment-id: o014.li2.chapter4.unbounded-derived-functors.i006
		\item Terdapat isomorfisme kanonik
		$\Hm^n\RHom(X,Y)\simeq
		\Hom_{\cate{D}(\mathcal{A})}(X,Y[n])=: \Ext^n(X,Y)$
		(Definisi~\ref{def:Ext-general}).
	\end{enumerate}
\end{teorema}
% segment-id: o014.li2.chapter4.unbounded-derived-functors.pf007
\begin{bukti}
	Argumennya hampir sama dengan Teorema~\ref{prop:RHom-bounded}. Sebagai
	contoh, anggap $\mathcal{A}$ mempunyai cukup banyak kompleks K-injektif.
	Kuncinya ialah menunjukkan bahwa
	$\cate{K}(\mathcal{A})^{\opp}\times\mathcal{I}$ bersifat
	$\Hom$-injektif.
% segment-id: o014.li2.chapter4.unbounded-derived-functors.l003
	\begin{itemize}
% segment-id: o014.li2.chapter4.unbounded-derived-functors.i007
		\item Tetapkan kompleks $X_1$. Funktor
		$(\cate{K}_\Pi\Hom)(X_1,\cdot)$ memetakan setiap kompleks K-injektif
		asiklik $X_2$ ke kompleks asiklik, sebab $X_2$ adalah $0$ dalam
		$\cate{K}(\mathcal{A})$ (Lema~\ref{prop:K-injective-zero}).
% segment-id: o014.li2.chapter4.unbounded-derived-functors.i008
		\item Tetapkan kompleks K-injektif $X_2$. Jika kompleks $X_1$ asiklik,
		definisi kompleks K-injektif menyiratkan bahwa
		$\Hom^\bullet\left(X_1,X_2\right)$ asiklik. Dengan menggunakan
		\eqref{eqn:CPiHom}, terlihat bahwa
		$(\cate{K}_\Pi\Hom)(X_1,X_2)$ juga asiklik.
	\end{itemize}
	Dengan demikian, syarat asiklik yang diperlukan untuk sifat
	$\Hom$-injektif telah diverifikasi, sedangkan syarat resolusinya langsung
	terpenuhi.
\end{bukti}

% segment-id: o014.li2.chapter4.unbounded-derived-functors.th002
\begin{teorema}\label{prop:RHom-adjoint}
	Tinjau pasangan funktor aditif adjoin antarkategori abelian
% segment-id: o014.li2.chapter4.unbounded-derived-functors.g001
	$\begin{tikzcd}
		F: \mathcal{A} \arrow[r, shift left] & \mathcal{A}' \arrow[l, shift left] : G
	\end{tikzcd}$.
	Misalkan $\mathcal{A}$ mempunyai cukup banyak kompleks K-projektif dan
	$\mathcal{A}'$ mempunyai cukup banyak kompleks K-injektif. Maka terdapat
	isomorfisme kanonik dalam $\cate{D}(\cate{Ab})$
% segment-id: o014.li2.chapter4.unbounded-derived-functors.q008
	\[ \RHom_{\mathcal{A'}}\left(\mathrm{L}F(X), Y\right) \simeq \RHom_{\mathcal{A}}\left(X, \mathrm{R}G(Y) \right), \]
	dengan $X\in\Obj\left(\cate{D}(\mathcal{A})\right)$ dan
	$Y\in\Obj\left(\cate{D}(\mathcal{A}')\right)$.
\end{teorema}
% segment-id: o014.li2.chapter4.unbounded-derived-functors.pf008
\begin{bukti}
	Argumennya sama dengan Teorema~\ref{prop:RHom-adjoint-bounded}. Satu-satunya
	perbedaan ialah bahwa resolusi injektif (atau projektif) diganti dengan
	resolusi K-injektif (atau K-projektif), dan
	Teorema~\ref{prop:K-resolution-homotopy} menggantikan
	Teorema~\ref{prop:resolution-homotopy} dalam pembahasan keunikan resolusi.
\end{bukti}

% segment-id: o014.li2.chapter4.unbounded-derived-functors.r001
\begin{catatan}
	Dengan menerapkan $\Hm^0$ pada kedua ruas dalam
	Teorema~\ref{prop:RHom-adjoint}, diperoleh adjungsi
% segment-id: o014.li2.chapter4.unbounded-derived-functors.q009
	\[ \Hom_{\cate{D}(\mathcal{A'})}\left(\mathrm{L}F(X), Y\right) \simeq \Hom_{\cate{D}(\mathcal{A})}\left(X, \mathrm{R}G(Y) \right). \]

	Karena funktor turunan yang diperoleh melalui resolusi selalu merupakan
	ekstensi Kan absolut dalam pengertian Definisi~\ref{def:absolute-Kan}
	(Proposisi~\ref{prop:localization-functor-abs}),
	Teorema~\ref{prop:Quillen-adjunction-gen} juga menyatakan bahwa
	$(\mathrm{L}F,\mathrm{R}G)$ merupakan pasangan adjoin. Keduanya memberikan
	isomorfisme adjungsi yang sama. Mengapa? Cukup ditinjau kasus ketika
	$X$ merupakan kompleks K-projektif dan $Y$ kompleks K-injektif. Dalam kasus
	ini, isomorfisme adjungsi dalam Teorema~\ref{prop:RHom-adjoint} diwujudkan
	secara konkret melalui unit $\eta$ dan kounit $\varepsilon$ pada tingkat
	$\cate{K}(\cdot)$ sebagai
% segment-id: o014.li2.chapter4.unbounded-derived-functors.q010
	\begin{align*}
		\Hom_{\cate{K}(\mathcal{A}')}\left( \cate{K}F(X), Y\right) \simeq \Hom_{\cate{K}(\mathcal{A})}\left( X, \cate{K}G(Y) \right).
	\end{align*}
	Di sisi lain, Teorema~\ref{prop:Quillen-adjunction-gen} menentukan
	$\underline{\eta}$ dan $\underline{\varepsilon}$ pada tingkat kategori
	turunan. Karakterisasi yang diberikan di sana cukup untuk menunjukkan bahwa
	keduanya kompatibel dengan $\eta$ dan $\varepsilon$. Perinciannya diserahkan
	kepada pembaca yang berminat.
\end{catatan}

% segment-id: o014.li2.chapter4.unbounded-derived-functors.p011
Sebagai penutup, sekumpulan syarat cukup lainnya untuk memperoleh funktor turunan
tak terbatas didasarkan pada keberhinggaan dimensi
(Definisi--Proposisi~\ref{def:dim-functor}). Meskipun teknik pembuktiannya tidak
melampaui cakupan buku ini, hasilnya akan dinyatakan tanpa bukti agar pembahasan
tidak menyimpang; lihat {\cite[Tags 07K6, 07K7]{stacks}}. Titik pangkalnya
adalah konstruksi funktor turunan terbatas dari
Teorema~\ref{prop:derived-via-resolution}.

% segment-id: o014.li2.chapter4.unbounded-derived-functors.pr003
\begin{proposisi}\label{prop:unbounded-derived-fd}
% segment-id: o014.li2.chapter4.unbounded-derived-functors.x004
	\index{dimensi}
	Misalkan $F:\mathcal{A}\to\mathcal{A}'$ funktor eksak kiri (atau eksak
	kanan) antarkategori abelian. Andaikan bahwa
% segment-id: o014.li2.chapter4.unbounded-derived-functors.l004
	\begin{compactitem}
% segment-id: o014.li2.chapter4.unbounded-derived-functors.i009
		\item relatif terhadap $F$, terdapat subkategori penuh aditif
		$\mathcal{A}^\flat$ dari $\mathcal{A}$ yang bertipe I (atau tipe P),
		seperti dalam Definisi--Proposisi~\ref{prop:F-injective-criterion};
% segment-id: o014.li2.chapter4.unbounded-derived-functors.i010
		\item funktor turunan terbatas yang bersesuaian,
		${}^+\mathrm{R}F$ (atau ${}^-\mathrm{L}F$), berdimensi hingga.
	\end{compactitem}
	Maka $\mathrm{R}F$ (atau $\cate{L}F$):
	$\cate{D}(\mathcal{A})\to\cate{D}(\mathcal{A}')$ ada. Bahkan,
	$\cate{K}(\mathcal{A}^\flat)$ membentuk subkategori $F$-injektif (atau
	$F$-projektif) dari $\cate{K}(\mathcal{A})$.
\end{proposisi}

% segment-id: o014.li2.chapter4.unbounded-derived-functors.ex001
\begin{contoh}\label{eg:Rlim-unbounded}
	Misalkan $\mathcal{A}$ mempunyai produk terhitung eksak dan cukup banyak
	objek injektif. Ambil funktor eksak kiri $F$ sebagai
	$\varprojlim:\cate{InvSys}(\mathcal{A})\to\mathcal{A}$.
	Teorema~\ref{prop:Rlim-as-holim} (iii) menunjukkan bahwa versi terbatas
	$\mathrm{R}\lim$ berdimensi paling besar $1$, sedangkan subkategori
	$\mathcal{R}\subset\cate{InvSys}(\mathcal{A})$ bertipe I relatif
	terhadapnya. Karena itu, hasil di atas memberikan funktor turunan tak terbatas
% segment-id: o014.li2.chapter4.unbounded-derived-functors.q011
	\[ \mathrm{R}\lim: \cate{D}(\cate{InvSys}(\mathcal{A})) \to \cate{D}(\mathcal{A}). \]

	Seperti dalam Teorema~\ref{prop:Rlim-as-holim} (ii), untuk setiap objek $X$
	dari $\cate{C}(\cate{InvSys}(\mathcal{A}))$ terdapat segitiga terbedakan
	dalam $\cate{D}(\mathcal{A})$
% segment-id: o014.li2.chapter4.unbounded-derived-functors.q012
	\[ \mathrm{R}\lim(X) \to \prod_k X_k \xrightarrow{\Delta_X} \prod_k X_k \xrightarrow{+1} . \]
	Secara khusus, $\mathrm{R}\lim(X)\simeq\holim(X)$. Berdasarkan
	Proposisi~\ref{prop:unbounded-derived-fd}, pembuktian semula berlaku tanpa
	perubahan lain: cukup ganti setiap $\cate{C}^+(\cdots)$ dengan
	$\cate{C}(\cdots)$.
\end{contoh}
